% !TeX root = ../main.tex
\chapter{Introdução}
\label{cap:introducao}

\section{Contexto e problema}

Ferramentas de análise estática de segurança (\textit{Static Application Security Testing}, SAST) produzem alertas que ajudam o desenvolvedor a localizar código que merece revisão. No caso da autorização, determinar a correção exige conhecer quais identidades podem realizar uma operação e em quais recursos. Essas regras podem estar distribuídas entre a configuração da aplicação, os metadados do endpoint e os serviços que decidem o acesso. Uma sugestão baseada apenas no alerta pode restringir uma operação a um papel inadequado ou permitir que uma identidade autenticada acesse um recurso para o qual ela não tem permissão.

O ASP.NET Core permite declarar políticas em controladores, ações e endpoints, além de avaliar a autorização explicitamente por meio de \texttt{IAuthorizationService}. Políticas podem combinar papéis, \textit{claims} e requisitos avaliados por \textit{handlers}. Quando a decisão depende da entidade carregada, como um documento associado a um proprietário, a verificação precisa levar esse recurso em conta \cite{aspnetpolicies,aspnetresource}. A correção precisa, portanto, considerar a operação, a política e as convenções do repositório.

O recorte adota a CWE-285 como categoria de referência e utiliza CWE-862 e CWE-863 para classificar os casos. A CWE-862 descreve a ausência da verificação exigida. A CWE-863 descreve uma verificação existente que decide o acesso de modo incorreto. A autorização por recurso pode aparecer em ambas as classes, conforme a natureza da falha observada \cite{cwe285,cwe862,cwe863}.

\section{Trabalhos relacionados}

A geração aumentada por recuperação (\textit{Retrieval-Augmented Generation}, RAG) acrescenta informações recuperadas à entrada do modelo de linguagem \cite{lewis2020rag}. Neste trabalho, essas informações são trechos de código e relações entre símbolos do repositório avaliado. O PatchAgent integra localização, geração e validação de correções \cite{yu2025patchagent}. O GitHub também oferece sugestões de correção em seu fluxo de alertas de análise de código \cite{githubautofix}. Esses trabalhos estabelecem precedentes para a remediação assistida. A questão deste projeto é medir o efeito da recuperação orientada ao modelo de autorização do ASP.NET Core, mantendo fixos o modelo de linguagem e o procedimento de avaliação.

A hipótese é que acrescentar contexto de autorização ao arquivo completo aumenta a proporção de correções que atendem simultaneamente aos três critérios, sem reduzir a proporção de correções seguras e funcionais. O contraste principal será entre o arquivo completo e o arquivo acrescido de contexto recuperado. A condição com apenas o método afetado permitirá observar quanto do ganho pode ser atribuído à leitura do arquivo.

\section{Objetivos}

\subsection{Objetivo geral}

Desenvolver e avaliar uma ferramenta de linha de comando que sugere correções de autorização em APIs ASP.NET Core com apoio de um modelo de linguagem e do contexto recuperado do repositório.

\subsection{Objetivos específicos}

\begin{enumerate}
    \item Normalizar alertas SARIF e manifestos experimentais e localizar os endpoints correspondentes com Roslyn.
    \item Recuperar as políticas, os requisitos, os \textit{handlers}, os serviços e as implementações relevantes para a operação afetada.
    \item Gerar propostas de correção em formato \textit{diff}, com justificativa, referências e indicação de incerteza.
    \item Validar as propostas em cópias isoladas e registrar evidências separadas de segurança, funcionalidade e aderência arquitetural.
    \item Comparar as condições de geração em um piloto reproduzível e avaliar a viabilidade de ampliar o conjunto de casos.
\end{enumerate}

\section{Escopo e estrutura do documento}

O protótipo abrange APIs ASP.NET Core em C\#, com controladores e Minimal APIs, e trata autorização ausente, política incorreta e falhas de autorização por recurso. A ferramenta atua na remediação após a indicação de um caso por um SAST ou pelo manifesto experimental. O escopo exclui a criação de um detector genérico, o suporte a outras linguagens ou classes de vulnerabilidade, a integração com IDE e a aplicação automática no repositório original.

O Capítulo~\ref{cap:detalhamento} descreve a arquitetura e os casos. O Capítulo~\ref{cap:cronograma} apresenta o cronograma. O Capítulo~\ref{cap:testes} define os testes e o experimento. O Capítulo~\ref{cap:riscos} analisa os riscos, e o Capítulo~\ref{cap:conclusao} discute a viabilidade do recorte.
